\begin{examplebox}
	\textbf{\textcolor{CExample}{例 1（基础）}}

	\textbf{\textcolor{CExample}{题目：}}
	已知球半径$R=5$，球心到截面的距离$d=3$，求截面圆的半径$r$。

	\textbf{\textcolor{CExample}{解题步骤：}}
	\begin{description}[style=nextline, leftmargin=2.8em, labelsep=0.5em, itemsep=0.45em]
		\item[\textbf{第一步（识别条件）：}] $R=5$，$d=3$，代入公式$r=\sqrt{R^{2}-d^{2}}$。
			\par\textbf{理由：} 公式直接给出截面半径与球半径和距离的关系。
		\item[\textbf{第二步（计算数值）：}] $r=\sqrt{5^{2}-3^{2}}=\sqrt{25-9}=\sqrt{16}=4$。
			\par\textbf{理由：} 按算术运算顺序计算平方、差、平方根。
		\item[\textbf{第三步（得出结论）：}] 截面圆的半径为$4$。
			\par\textbf{理由：} 计算无误，结果为正。
	\end{description}

	\textbf{\textcolor{CExample}{关键结论：}}
	\textbf{$r=4$。}

	\medskip
	\textbf{\textcolor{CExample}{例 2（稍变形）}}

	\textbf{\textcolor{CExample}{题目：}}
	已知球半径$R=10$，截面圆半径$r=6$，求球心到平面的距离$d$。

	\textbf{\textcolor{CExample}{解题步骤：}}
	\begin{description}[style=nextline, leftmargin=2.8em, labelsep=0.5em, itemsep=0.45em]
		\item[\textbf{第一步（选用逆公式）：}] 使用逆公式$d=\sqrt{R^{2}-r^{2}}$。
			\par\textbf{理由：} 已知$R$和$r$，反求$d$。
		\item[\textbf{第二步（代入计算）：}] $d=\sqrt{10^{2}-6^{2}}=\sqrt{100-36}=\sqrt{64}=8$。
			\par\textbf{理由：} 算术计算。
		\item[\textbf{第三步（写出答案）：}] 球心到截面的距离为$8$。
			\par\textbf{理由：} 结果合理，$d<R$满足条件。
	\end{description}

	\textbf{\textcolor{CExample}{关键结论：}}
	\textbf{$d=8$。}

	\medskip
	\textbf{\textcolor{CExample}{例 3（综合应用）}}

	\textbf{\textcolor{CExample}{题目：}}
	一个半径为$13$的球，被一平面截得的截面面积为$25\pi$，求球心到截面的距离$d$。

	\textbf{\textcolor{CExample}{解题步骤：}}
	\begin{description}[style=nextline, leftmargin=2.8em, labelsep=0.5em, itemsep=0.45em]
		\item[\textbf{第一步（由面积求半径）：}] 由$S=\pi r^{2}=25\pi$，得$r^{2}=25$，$r=5$（取正）。
			\par\textbf{理由：} 截面是圆，圆面积公式$S=\pi r^{2}$。
		\item[\textbf{第二步（利用公式求距离）：}] $d=\sqrt{R^{2}-r^{2}}=\sqrt{13^{2}-5^{2}}=\sqrt{169-25}=\sqrt{144}=12$。
			\par\textbf{理由：} 球截面半径公式的逆用，勾股定理。
		\item[\textbf{第三步（得出结论）：}] 球心到截面的距离为$12$。
			\par\textbf{理由：} 距离在$0\le d<R$范围内，结果有效。
	\end{description}

	\textbf{\textcolor{CExample}{关键结论：}}
	\textbf{$d=12$。}
\end{examplebox}
